% ============================================================
% CAPITOLO 11 - ARMONIZZARE LA SCALA: REGOLA DELL'OTTAVA PRELIMINARE
% ============================================================
\chapter{Armonizzare la scala: la Regola dell'Ottava preliminare}
\chaptermark{La Regola dell'Ottava preliminare}
\label{chap:ottava_preliminare}

La \textbf{Regola dell'Ottava} è uno schema pratico che indica come armonizzare ogni grado della scala quando compare al basso. Ha origine nella pratica del basso continuo barocco ed è soprattutto uno strumento didattico per capire la logica dell'armonizzazione.

\section{La regola dell'ottava con sole triadi}
\label{sec:regola_ottava_funzioni}

\begin{definizione}[Le tre versioni della Regola dell'Ottava]
	\label{def:versioni_regola_ottava}
	Il manuale presenta la Regola dell'Ottava in \textbf{tre versioni progressive}:
	\begin{enumerate}
		\item \textbf{Versione preliminare (questo capitolo):} usa soltanto \textbf{triadi}, in stato fondamentale e rivoltate.
		\item \textbf{Versione intermedia (Capitolo \ref{chap:settima_dominante}):} aggiunge la \textbf{settima di dominante} e i suoi rivolti.
		\item \textbf{Versione completa (Capitolo \ref{chap:regola_ottava_completa}):} include tutte le settime diatoniche.
	\end{enumerate}
	Procedere per gradi evita di affrontare subito tutte le possibilità armoniche. La parte cromatica della regola è nel Capitolo \ref{chap:regola_ottava_cromatica}.
\end{definizione}

Con le sole triadi la scala si armonizza alternando stati fondamentali e primi rivolti. Il rivolto mantiene il basso in movimento melodico senza appesantire la struttura con troppi stati fondamentali.

\section{Movimenti del basso e assegnazione delle funzioni}
\label{subsec:principio_regola_ottava}

Quando il basso procede per grado congiunto, l'accordo da usare dipende dalla posizione della nota nella scala. Le sei regole seguenti danno lo schema per i due modi, in senso ascendente e discendente. Nel minore armonico la sensibile è alzata; nel minore melodico la scala ascendente alza anche il VI grado, mentre quella discendente coincide con il minore naturale.

\begin{regola}[Regola dell'Ottava (triadi): scala maggiore ascendente]
	\label{reg:regola_ottava_maggiore_asc}
	\begin{itemize}
		\item \textbf{I grado:} accordo di I in stato fondamentale (5)
		\item \textbf{II grado:} accordo di VII in primo rivolto (6)
		\item \textbf{III grado:} accordo di I in primo rivolto (6)
		\item \textbf{IV grado:} accordo di II in primo rivolto (6)
		\item \textbf{V grado:} accordo di V in stato fondamentale (5), oppure \sigla{I}{6}{4} di passaggio
		\item \textbf{VI grado:} accordo di IV in primo rivolto (6)
		\item \textbf{VII grado:} accordo di V in primo rivolto (6)
		\item \textbf{VIII grado:} accordo di I in stato fondamentale (5)
	\end{itemize}
\end{regola}

\begin{regola}[Regola dell'Ottava (triadi): scala maggiore discendente]
	\label{reg:regola_ottava_maggiore_disc}
	\begin{itemize}
		\item \textbf{VIII grado:} accordo di I in stato fondamentale (5)
		\item \textbf{VII grado:} accordo di V in primo rivolto (6)
		\item \textbf{VI grado:} accordo di IV in primo rivolto (6)
		\item \textbf{V grado:} accordo di V in stato fondamentale (5)
		\item \textbf{IV grado:} accordo di II in primo rivolto (6)
		\item \textbf{III grado:} accordo di I in primo rivolto (6)
		\item \textbf{II grado:} accordo di VII in primo rivolto (6)
		\item \textbf{I grado:} accordo di I in stato fondamentale (5)
	\end{itemize}
\end{regola}

\begin{regola}[Regola dell'Ottava (triadi): scala minore armonica ascendente]
	\label{reg:regola_ottava_minore_asc}
	\begin{itemize}
		\item \textbf{I grado:} accordo di I in stato fondamentale (5)
		\item \textbf{II grado:} accordo di VII$^\circ$ in primo rivolto (6)
		\item \textbf{III grado:} accordo di I in primo rivolto (6)
		\item \textbf{IV grado:} accordo di II$^\circ$ in primo rivolto (6)
		\item \textbf{V grado:} accordo di V in stato fondamentale (5), oppure \sigla{I}{6}{4} di passaggio
		\item \textbf{VI grado:} accordo di IV in primo rivolto (6)
		\item \textbf{VII grado:} accordo di V in primo rivolto (6)
		\item \textbf{VIII grado:} accordo di I in stato fondamentale (5)
	\end{itemize}
	Fra il VI e il VII grado il basso forma una seconda aumentata (Fa-Sol$\sharp$ in La minore): per questo, salendo, si preferisce di solito la scala minore melodica.
\end{regola}

\begin{regola}[Regola dell'Ottava (triadi): scala minore armonica discendente]
	\label{reg:regola_ottava_minore_disc}
	\begin{itemize}
		\item \textbf{VIII grado:} accordo di I in stato fondamentale (5)
		\item \textbf{VII grado:} accordo di V in primo rivolto (6)
		\item \textbf{VI grado:} accordo di IV in primo rivolto (6)
		\item \textbf{V grado:} accordo di V in stato fondamentale (5)
		\item \textbf{IV grado:} accordo di II$^\circ$ in primo rivolto (6)
		\item \textbf{III grado:} accordo di I in primo rivolto (6)
		\item \textbf{II grado:} accordo di VII$^\circ$ in primo rivolto (6)
		\item \textbf{I grado:} accordo di I in stato fondamentale (5)
	\end{itemize}
\end{regola}

\begin{regola}[Regola dell'Ottava (triadi): scala minore melodica ascendente]
	\label{reg:regola_ottava_melodica_asc}
	Con il VI e il VII grado alzati, la triade del II grado diventa minore (Si-Re-Fa$\sharp$ in La minore) e quella del IV grado maggiore (Re-Fa$\sharp$-La), come nella scala maggiore.
	\begin{itemize}
		\item \textbf{I grado:} accordo di I in stato fondamentale (5)
		\item \textbf{II grado:} accordo di VII$^\circ$ in primo rivolto (6)
		\item \textbf{III grado:} accordo di I in primo rivolto (6)
		\item \textbf{IV grado:} accordo di II in primo rivolto (6)
		\item \textbf{V grado:} accordo di V in stato fondamentale (5), oppure \sigla{I}{6}{4} di passaggio
		\item \textbf{VI grado:} accordo di IV in primo rivolto (6)
		\item \textbf{VII grado:} accordo di V in primo rivolto (6)
		\item \textbf{VIII grado:} accordo di I in stato fondamentale (5)
	\end{itemize}
\end{regola}

\begin{regola}[Regola dell'Ottava (triadi): scala minore melodica discendente]
	\label{reg:regola_ottava_melodica_disc}
	Si usano i gradi naturali (VI e VII abbassati), come nella scala minore naturale.
	\begin{itemize}
		\item \textbf{VIII grado:} accordo di I in stato fondamentale (5)
		\item \textbf{VII grado:} accordo di V minore in primo rivolto (6), con il VII grado naturale al basso (Mi-Sol-Si su Sol in La minore)
		\item \textbf{VI grado:} accordo di IV in primo rivolto (6)
		\item \textbf{V grado:} accordo di V in stato fondamentale (5)
		\item \textbf{IV grado:} accordo di II$^\circ$ in primo rivolto (6)
		\item \textbf{III grado:} accordo di I in primo rivolto (6)
		\item \textbf{II grado:} accordo di VII$^\circ$ in primo rivolto (6)
		\item \textbf{I grado:} accordo di I in stato fondamentale (5)
	\end{itemize}
\end{regola}

\section{Basso per salto o per grado congiunto}
\label{sec:basso_salti_gradi}

Il movimento melodico (grado congiunto) e quello armonico (salto) del basso decidono lo stato dell'accordo.

\begin{regola}[Scelta dello stato in base al movimento del basso]
	\label{reg:stato_basso}
	\begin{itemize}
		\item \textbf{Basso per grado congiunto:} si applica la Regola dell'Ottava.
		\item \textbf{Basso che salta:} si preferisce lo \textbf{stato fondamentale}, perché il salto segnala un cambio di pilastro armonico.
		\item \textbf{Eccezione del V grado:} la dominante può sempre essere usata in stato fondamentale, per sottolineare la forza delle cadenze, qualunque sia il movimento del basso.
	\end{itemize}
\end{regola}

\begin{regola}[Salto d'ottava al basso]
	\label{reg:salto_ottava_basso}
	Quando il basso compie un salto d'ottava, in su o in giù, la seconda nota del salto si numera di norma con il \textbf{6} (primo rivolto della triade), per avere un arpeggio armonico più fluido e vario.
\end{regola}

\section{Il V grado: dominante o \sigla{I}{6}{4}}
\label{sec:secondo_rivolto_regola}

Sul V grado la regola ammette due letture. Lo stato fondamentale della dominante è la scelta più comune e dà stabilità; si può usare anche il secondo rivolto della tonica, \sigla{I}{6}{4}, in due modi:

\begin{itemize}
	\item \textbf{\sigla{I}{6}{4} di passaggio:} il basso tocca il V grado muovendosi tra il IV e il VI (o viceversa), con un collegamento morbido.
	\item \textbf{\sigla{I}{6}{4} cadenzale:} il basso sul V grado viene «appoggiato» da quarta e sesta, che poi risolvono sulla terza e sulla quinta della dominante (Sezione \ref{sec:cadenziale}).
\end{itemize}

Perché il \sigla{I}{6}{4} sul V grado suoni «sospeso» è spiegato nell'Approfondimento \ref{phys:i64_sospeso}.
